\section{Block Lyapunov Identities for the Calibrated Clamp}
\label{app:block_identities}

This appendix isolates the algebra used in the paper's shorting argument.  The
only inputs are the block-triangular calibrated operator family, the trace-class
stationary covariance, and the Cameron--Martin finite-energy assumption on the
outgoing drift channel.

\subsection{Block identities}

Write $\mathcal{H}=\mathcal{H}_k\oplus\mathcal{H}_I$ with
$\mathcal{H}_k=\mathbb{R}\phi_k$.  Under the calibrated clamp,
\begin{equation}
 \mathcal{A}^{(\lambda)}=
 \begin{pmatrix}-\lambda&0\\a_{Ik}&\mathcal{A}_{II}\end{pmatrix},
 \qquad
 \mathcal{D}^{(\lambda)}=
 \begin{pmatrix}\varepsilon_0/\lambda&0\\0&D_{II}\end{pmatrix}.
 \label{eq:exact_blocks}
\end{equation}
Let $\Sigma_\lambda$ be the stationary covariance solving
\[
  \mathcal{A}^{(\lambda)}\Sigma_\lambda
  +\Sigma_\lambda(\mathcal{A}^{(\lambda)})^*
  +\mathcal{D}^{(\lambda)}=0.
\]

\begin{lemma}[Block Lyapunov identities]
\label{lem:block_lyapunov_identities}
For every $\lambda>0$,
\begin{equation}
  \Sigma_{\lambda,kk}=\frac{\varepsilon_0}{2\lambda^2},
  \label{eq:Sigmakk_exact}
\end{equation}
and the target-to-bulk cross covariance is
\begin{equation}
  \Sigma_{\lambda,Ik}
  =\frac{\varepsilon_0}{2\lambda^2}
    (\lambda I-\mathcal{A}_{II})^{-1}a_{Ik}.
  \label{eq:SigmaIk_exact_vector}
\end{equation}
\end{lemma}

\begin{proof}
The $kk$ block of the Lyapunov equation is
\[
  -2\lambda\Sigma_{\lambda,kk}+\frac{\varepsilon_0}{\lambda}=0,
\]
which gives \eqref{eq:Sigmakk_exact}.  The $Ik$ block is
\[
  (\mathcal{A}_{II}-\lambda I)\Sigma_{\lambda,Ik}
  +a_{Ik}\Sigma_{\lambda,kk}=0.
\]
Substituting \eqref{eq:Sigmakk_exact} gives
\eqref{eq:SigmaIk_exact_vector}.
\end{proof}

\subsection{Cameron--Martin energy}

Let $Q_0\coloneqq\Sigma_{II}^{(0)}$ be the background bulk covariance and
$H_{Q_0}\coloneqq\Range(Q_0^{1/2})$ with energy norm
$|u|_{Q_0}=\|Q_0^{-1/2}u\|_{\mathcal{H}_I}$.  Assumption~\ref{assum:cm_stability}
states that $a_{Ik}\in H_{Q_0}$ and that the bulk resolvent is bounded on this
energy space:
\begin{equation}
  |(\lambda I-\mathcal{A}_{II})^{-1}u|_{Q_0}
  \le \frac{M}{\lambda+\omega}|u|_{Q_0}.
  \label{eq:cm_resolvent}
\end{equation}

\begin{lemma}[Finite-energy leakage]
\label{lem:finite_energy_leakage}
The cross covariance belongs to $H_{Q_0}$ and satisfies
\begin{equation}
 |\Sigma_{\lambda,Ik}|_{Q_0}
 \le
 \frac{\varepsilon_0M}{2\lambda^2(\lambda+\omega)}
 |a_{Ik}|_{Q_0}
 =O(\lambda^{-3}).
 \label{eq:cross_energy}
\end{equation}
\end{lemma}

\begin{proof}
Apply \eqref{eq:cm_resolvent} to
\eqref{eq:SigmaIk_exact_vector}.  This proves both the range statement and the
energy bound.
\end{proof}

\section{Anderson--Trapp Shorting and the Rank-One Residual}
\label{app:shorting_residual}

In finite dimensions, the Schur complement is only the bounded-inverse
shadow of Anderson--Trapp shorting.  The Hilbert space argument therefore
starts from shorting and recovers the Schur expression only when the
complementary block is boundedly invertible.

\begin{lemma}[Schur complement as the bounded-inverse shadow of shorting]
\label{lem:schur_shadow_shorting}
Let $\mathcal{E}$ and $\mathcal{F}$ be Hilbert spaces and let
\[
  T=\begin{pmatrix}A&B^*\\B&D\end{pmatrix}\ge0
\]
act on $\mathcal{E}\oplus\mathcal{F}$.  If $D$ is boundedly invertible on
$\mathcal{F}$, then the Anderson--Trapp short of $T$ to $\mathcal{E}$ has the
quadratic form
\begin{equation}
  \langle x,\mathcal{S}_{\mathcal{E}}(T)x\rangle
  =
  \langle x,Ax\rangle-\langle Bx,D^{-1}Bx\rangle.
  \label{eq:short_shadow}
\end{equation}
Thus $B^*D^{-1}B$ is exactly the variational shorting subtraction in the special
case where $D^{-1}$ is a bounded operator.
\end{lemma}

\begin{proof}
By the variational definition of the shorted operator,
\[
  \langle x,\mathcal{S}_{\mathcal{E}}(T)x\rangle
  =
  \inf_{y\in\mathcal{F}}
  \left\langle
  \begin{pmatrix}x\\y\end{pmatrix},
  T
  \begin{pmatrix}x\\y\end{pmatrix}
  \right\rangle .
\]
The displayed quadratic form equals
$\langle x,Ax\rangle+2\operatorname{Re}\langle Bx,y\rangle+\langle y,Dy\rangle$.
Since $D$ is boundedly invertible and positive, the unique minimizer is
$y=-D^{-1}Bx$.  Substitution gives \eqref{eq:short_shadow}.
\end{proof}

For the covariance block $\Sigma_\lambda$ on
$\mathbb{C}\phi_k\oplus\mathcal{H}_I$, the shorting subtraction is the
variational quantity
\begin{equation}
 R_\lambda
 =
 \sup_{\langle u,\Sigma_{\lambda,II}u\rangle>0}
 \frac{|\langle u,\Sigma_{\lambda,Ik}\rangle|^2}
      {\langle u,\Sigma_{\lambda,II}u\rangle}.
 \label{eq:schur_variational}
\end{equation}
The null directions of the denominator do not change the supremum for a
non-negative block operator.

\begin{proposition}[Rank-one shorted residual]
\label{prop:rank_one_shorted_residual}
Under Assumptions~\ref{assum:linear_second_moment} and
\ref{assum:cm_stability},
\begin{equation}
  S_{k,\lambda}
  =\Sigma_{\lambda,kk}-R_\lambda,
  \qquad
  0\le R_\lambda\le |\Sigma_{\lambda,Ik}|_{Q_0}^2=O(\lambda^{-6}).
  \label{eq:shorted_residual_bound}
\end{equation}
Consequently
\begin{equation}
  S_{k,\lambda}
  =
  \frac{\varepsilon_0}{2\lambda^2}-O(\lambda^{-6}),
  \qquad
  m_2^{\mathrm{Schur}}=\frac{\varepsilon_0}{2},
  \label{eq:Schur_asympt}
\end{equation}
and
\[
  S_{k,\lambda}^{-1}
  =
  \frac{2}{\varepsilon_0}\lambda^2+O(\lambda^{-2}).
\]
\end{proposition}

\begin{proof}
The first identity is the scalar form of Anderson--Trapp shorting applied to
the target coordinate.  Because the target noise and the free noise are
independent, the bulk covariance dominates the background covariance:
$\Sigma_{\lambda,II}=Q_0+\operatorname{Cov}(Y_I)\succeq Q_0$.  Douglas range
inclusion gives the contraction of energy norms
\[
  \|\Sigma_{\lambda,II}^{-1/2}v\|
  \le \|Q_0^{-1/2}v\|
  \quad (v\in H_{Q_0}).
\]
Using this inequality in \eqref{eq:schur_variational} and then applying
\eqref{eq:cross_energy} yields \eqref{eq:shorted_residual_bound}.  The
asymptotic formula follows from \eqref{eq:Sigmakk_exact}, and the reciprocal
expansion follows by expanding
$S_{k,\lambda}^{-1}$ around $\varepsilon_0/(2\lambda^2)$.
\end{proof}

\section{Galerkin Approximation and the Failure of Direct Inversion}
\label{app:galerkin_shorting}

Galerkin approximation is used here only to approximate trace-class covariances
and nested variational shorting problems.  It is not a license to replace the
Hilbert-space short by a global inverse of a compact covariance block.

\begin{lemma}[Trace-norm convergence of Galerkin covariances]
\label{lem:galerkin_trace_norm}
Let $\Pi_n:\mathcal{H}\to\mathcal{H}_n$ be finite-rank orthogonal projections
with $\Pi_n\phi_k=\phi_k$, $\Pi_n\mathcal{H}_I\subset\mathcal{H}_I$, and
$\Pi_n\to I$ strongly.  For fixed $\lambda>0$, set
\[
  \mathcal{A}_n^{(\lambda)}
  =\Pi_n\mathcal{A}^{(\lambda)}|_{\mathcal{H}_n},
  \qquad
  \mathcal{D}_n^{(\lambda)}
  =\Pi_n\mathcal{D}^{(\lambda)}\Pi_n .
\]
Assume that $e^{t\mathcal{A}_n^{(\lambda)}}x\to
e^{t\mathcal{A}^{(\lambda)}}x$ for every $x$, uniformly for $t$ in compact
intervals, that the semigroups are uniformly exponentially stable, satisfying
\[
  \|e^{t\mathcal{A}_n^{(\lambda)}}\|, \
  \|e^{t\mathcal{A}^{(\lambda)}}\|\le Me^{-\beta t},
\]
and that $\mathcal{D}_n^{(\lambda)}\to\mathcal{D}^{(\lambda)}$ in
$\mathcal{S}_1(\mathcal{H})$.  Then
\[
  \Sigma_{n,\lambda}
  =
  \int_0^\infty
  e^{t\mathcal{A}_n^{(\lambda)}}
  \mathcal{D}_n^{(\lambda)}
  e^{t\mathcal{A}_n^{(\lambda)*}}\,dt
  \longrightarrow
  \Sigma_\lambda
  \quad\text{in }\mathcal{S}_1(\mathcal{H}).
\]
\end{lemma}

\begin{proof}
Omitting the superscript $(\lambda)$, write the difference as a Bochner
integral and split the integrand:
\begin{multline*}
  \|e^{t\mathcal{A}_n}\mathcal{D}_ne^{t\mathcal{A}_n^*}
    -e^{t\mathcal{A}}\mathcal{D}e^{t\mathcal{A}^*}\|_{\mathcal{S}_1}
  \le
  M^2e^{-2\beta t}\|\mathcal{D}_n-\mathcal{D}\|_{\mathcal{S}_1}\\
  +
  \|e^{t\mathcal{A}_n}\mathcal{D}e^{t\mathcal{A}_n^*}
    -e^{t\mathcal{A}}\mathcal{D}e^{t\mathcal{A}^*}\|_{\mathcal{S}_1}.
\end{multline*}
The first term tends to zero.  For the second, Gr\"umm's theorem for trace
ideals implies trace-norm convergence from strong convergence of the left and
right semigroups, uniform boundedness, and $\mathcal{D}\in\mathcal{S}_1$
\cite{simon2005trace}.  The full integrand is dominated by
\[
  M^2e^{-2\beta t}
  \bigl(\|\mathcal{D}_n\|_{\mathcal{S}_1}
        +\|\mathcal{D}\|_{\mathcal{S}_1}\bigr),
\]
with a finite uniform trace-norm bound because
$\mathcal{D}_n\to\mathcal{D}$ in $\mathcal{S}_1$.  Dominated convergence for
Bochner integrals gives the claim.
\end{proof}

\begin{lemma}[Galerkin convergence of variational shorting]
\label{lem:galerkin_variational_shorting}
For every block-preserving Galerkin truncation satisfying the hypotheses of
\Cref{lem:galerkin_trace_norm} and the uniform finite-energy condition
\[
  a_{Ik}^{(n)}\in H_{Q_0^{(n)}},
  \qquad
  \sup_n |a_{Ik}^{(n)}|_{Q_0^{(n)}}<\infty,
\]
the truncated residual satisfies
\[
  S_{k,\lambda}^{(n)}
  =
  \frac{\varepsilon_0}{2\lambda^2}-R_{n,\lambda},
  \qquad
  0\le R_{n,\lambda}\le C\lambda^{-6},
\]
with $C$ independent of $n$.  Hence
$m_2^{\mathrm{Schur},(n)}=\varepsilon_0/2$ for every admissible truncation.
Moreover, for fixed $\lambda$, nested finite-dimensional test spaces in the
bulk Cameron--Martin space give monotone Rayleigh--Ritz convergence of the
variational subtraction to $R_\lambda$.
\end{lemma}

\begin{proof}
The block-preserving truncation has the same triangular form as
\eqref{eq:exact_blocks}, so the finite-dimensional Lyapunov blocks give
\[
  \Sigma_{n,\lambda,kk}=\frac{\varepsilon_0}{2\lambda^2},
  \qquad
  \Sigma_{n,\lambda,Ik}
  =
  \frac{\varepsilon_0}{2\lambda^2}
  (\lambda I-\mathcal{A}_{II}^{(n)})^{-1}a_{Ik}^{(n)}.
\]
The uniform finite-energy assumption gives
$|\Sigma_{n,\lambda,Ik}|_{Q_0^{(n)}}=O(\lambda^{-3})$ uniformly in $n$.
The same Douglas domination argument used in
\Cref{prop:rank_one_shorted_residual} then yields
$0\le R_{n,\lambda}\le C\lambda^{-6}$ and the stated value of
$m_2^{\mathrm{Schur},(n)}$.

For the variational convergence statement, let
$C_\lambda=\Sigma_{\lambda,II}$ and let $V_m$ be nested finite-dimensional
subspaces whose union is dense in $H_{C_\lambda}=\Range(C_\lambda^{1/2})$.
Then
\[
  R_\lambda^{[m]}
  =
  \sup_{\substack{u\in V_m\\ \langle u,C_\lambda u\rangle>0}}
  \frac{|\langle u,\Sigma_{\lambda,Ik}\rangle|^2}
       {\langle u,C_\lambda u\rangle}
\]
is monotone increasing in $m$ and bounded above by $R_\lambda$.  Density in the
Cameron--Martin energy norm gives $R_\lambda^{[m]}\uparrow R_\lambda$.  In an
eigenbasis of $C_\lambda$, this is the Rayleigh--Ritz identity
\[
  R_\lambda^{[m]}
  =
  \sum_{j=1}^m
  \frac{|\langle\Sigma_{\lambda,Ik},e_j\rangle|^2}{\sigma_j}
  \uparrow
  \|\Sigma_{\lambda,Ik}\|_{H_{C_\lambda}}^2.
\]
\end{proof}

The obstruction to direct inversion is not numerical bookkeeping.  In infinite
dimension the compact covariance block has no bounded inverse on the ambient
Hilbert space.  In finite-dimensional Galerkin matrices the condition number
$\kappa(D_N)$ diverges as the smallest covariance eigenvalue tends to zero, so
the formula $B_N^*D_N^{-1}B_N$ amplifies truncation and roundoff errors.  Even
with exact arithmetic, the limits do not commute: fixing $N$ and then sending
$\lambda\to\infty$ resolves a different finite matrix problem than first
passing to the trace-class covariance and then taking the calibrated shorting
limit.  The stable object is therefore the nested variational short, not direct
matrix inversion.

\section{Gaussian Radon Conditioning and the Fredholm Finite Part}
\label{app:gaussian_fredholm}

The Gaussian and Fredholm layer is the Radon representation of the shorted
residual constructed above.  It is not a second mechanism: disintegration
supplies the exact-evidence law, while the shorted covariance residual supplies
the finite conditional variance used by the Gaussian entropy calculation.

\subsection{Pointwise disintegration}

Let $\mu_\lambda=\mathcal{N}(m_\lambda,\Sigma_\lambda)$ be a Radon Gaussian
measure on $\mathcal{H}=\mathcal{H}_k\oplus\mathcal{H}_I$, and let
$\pi_k(X)=\langle X,\phi_k\rangle$.

\begin{lemma}[Continuous disintegration and exact target zeros]
\label{lem:disintegration}
For each $t\in\mathbb{R}$, there is a unique weakly continuous Gaussian
disintegration member
$\mu_\lambda^t=\mathcal{N}(m_\lambda^t,\Sigma_\lambda^t)$ supported on
$X_k=t$, where
\[
  m_\lambda^t=
  \begin{pmatrix}
    t\\
    m_{\lambda,I}
    +\Sigma_{\lambda,Ik}\Sigma_{\lambda,kk}^{-1}(t-m_{\lambda,k})
  \end{pmatrix},
\]
and
\[
  \Sigma_\lambda^t=
  \begin{pmatrix}
    0&0\\
    0&\Sigma_{\lambda,II}
      -\Sigma_{\lambda,Ik}\Sigma_{\lambda,kk}^{-1}\Sigma_{\lambda,kI}
  \end{pmatrix}.
\]
At $t=x_k^*$, the observed law $\mu_{\mathrm{obs}}=\mu_\lambda^{x_k^*}$
satisfies
\[
  \mathbb{E}_{\mathrm{obs}}[X_k]-x_k^*=0,
  \qquad
  \operatorname{Var}_{\mathrm{obs}}(X_k)=0.
\]
\end{lemma}

\begin{proof}
Regular conditional probabilities exist because $\mathcal{H}$ is Polish.  The
displayed family is the weakly continuous Gaussian disintegration along the
continuous scalar observation map, and this continuous version is unique
\cite{bogachev1998gaussian,lagatta2013continuous}.  Evaluating at $x_k^*$ fixes
the target coordinate almost surely.
\end{proof}

\subsection{Gaussian energy projection and entropy}

Let $C_\lambda=\Sigma_{\lambda,II}$.  Since $C_\lambda$ is trace-class, its
inverse is not a bounded ambient operator; the relevant object is the
Cameron--Martin space $H_{C_\lambda}=\Range(C_\lambda^{1/2})$.

\begin{lemma}[Measurable Cameron--Martin energy projection]
\label{lem:conditional_predictor}
The cross covariance belongs to $H_{C_\lambda}$ and
\[
  \|\Sigma_{\lambda,Ik}\|_{H_{C_\lambda}}
  \le |\Sigma_{\lambda,Ik}|_{Q_0}=O(\lambda^{-3}).
\]
The conditional mean projection is the measurable Wiener functional
\[
  \mu_{k|I}(X_I)
  =
  m_{\lambda,k}
  +W_{\Sigma_{\lambda,Ik}}(X_I-m_{\lambda,I}),
\]
and, for all sufficiently large $\lambda$,
\[
  \operatorname{Var}(X_k\mid X_I)
  =
  \Sigma_{\lambda,kk}
  -\|\Sigma_{\lambda,Ik}\|_{H_{C_\lambda}}^2
  =S_{k,\lambda}>0.
\]
\end{lemma}

\begin{proof}
The domination $\Sigma_{\lambda,II}\succeq Q_0$ gives
$H_{Q_0}\subseteq H_{C_\lambda}$ by Douglas range inclusion
\cite{douglas1966majorization}.  Combining this inclusion with
\eqref{eq:cross_energy} gives the norm bound.  The range condition is exactly
the condition for the measurable Wiener functional associated with
$\Sigma_{\lambda,Ik}$ to exist under the bulk Gaussian law
\cite{bogachev1998gaussian}.  Positivity follows from
\Cref{prop:rank_one_shorted_residual}.
\end{proof}

\begin{lemma}[Feldman--Hajek equivalence and trace-anomaly decay]
\label{lem:feldman_hajek}
There is $\lambda_0>0$ such that for $\lambda\ge\lambda_0$ the bulk marginals
$\mu_{\mathrm{obs},I}$ and $\mu_{\lambda,I}$ are equivalent and
\[
  2\KL(\mu_{\mathrm{obs},I}\parallel\mu_{\lambda,I})
  =
  \frac{(x_k^*-m_{\lambda,k})^2}{\Sigma_{\lambda,kk}}\rho_\lambda
  -\log(1-\rho_\lambda)-\rho_\lambda,
\]
where
\[
  \rho_\lambda
  =
  \frac{\|\Sigma_{\lambda,Ik}\|_{H_{C_\lambda}}^2}
       {\Sigma_{\lambda,kk}}
  =
  1-\frac{S_{k,\lambda}}{\Sigma_{\lambda,kk}}
  =O(\lambda^{-4}).
\]
In particular,
$\KL(\mu_{\mathrm{obs},I}\parallel\mu_{\lambda,I})=O(\lambda^{-4})$.
\end{lemma}

\begin{proof}
The covariance change under disintegration is the rank-one operator
$\Sigma_{\lambda,Ik}\Sigma_{\lambda,kk}^{-1}\Sigma_{\lambda,kI}$.  By
\Cref{lem:conditional_predictor}, its relative covariance perturbation has the
single non-unit eigenvalue $1-\rho_\lambda$.  For large $\lambda$,
$\rho_\lambda<1$, so Feldman--Hajek equivalence applies
\cite{feldman1958equivalence,bogachev1998gaussian}.  The Gaussian relative
entropy formula gives the displayed expression.  Assumption~\ref{assum:gaussian_clamp}
has $m_{\lambda,k}-x_k^*=O(\lambda^{-1})$, while
$\Sigma_{\lambda,kk}=O(\lambda^{-2})$ and $\rho_\lambda=O(\lambda^{-4})$; the
logarithmic remainder is $O(\rho_\lambda^2)$.  Hence the entropy is
$O(\lambda^{-4})$.
\end{proof}

\subsection{Mollified divergence and finite part}

For $\delta>0$, set
\[
  \nu_{\lambda,\delta}
  =
  \mu_{\mathrm{obs},I}\otimes\mathcal{N}(x_k^*,\delta^2)
\]
under the product identification
$\mathcal{H}=\mathcal{H}_I\times\mathcal{H}_k$.

\begin{proposition}[Mollification limit]
\label{prop:mollification_limit}
For $\lambda\ge\lambda_0$ as in \Cref{lem:feldman_hajek}, define
\[
  C_{\lambda,\delta}
  =
  \frac12\left(\log\frac{S_{k,\lambda}}{\delta^2}-1\right).
\]
Then
\[
  \lim_{\delta\downarrow0}
  \left[
    \KL(\nu_{\lambda,\delta}\parallel\mu_\lambda)
    -C_{\lambda,\delta}
  \right]
  =
  \mathfrak{J}_\lambda(\mu_{\mathrm{obs}};\mu_\lambda).
\]
\end{proposition}

\begin{proof}
Let $r_\lambda(X)=X_k-\mu_{k|I}(X_I)$.  The entropy chain rule and the scalar
Gaussian formula give
\[
  \KL(\nu_{\lambda,\delta}\parallel\mu_\lambda)
  =
  \KL(\mu_{\mathrm{obs},I}\parallel\mu_{\lambda,I})
  +C_{\lambda,\delta}
  +\frac{\delta^2}{2S_{k,\lambda}}
  +\frac{1}{2S_{k,\lambda}}
   \mathbb{E}_{\mu_{\mathrm{obs}}}[r_\lambda(X)^2].
    \]
    The $\delta^2$ term vanishes as $\delta\downarrow0$, and the remaining terms
    are exactly Definition~\ref{def:unified_action}.
\end{proof}

\begin{lemma}[Weak scheme-independence]
\label{lem:weak_scheme_independence}
For an evidence level $t\in\mathbb{R}$, let $\mu_\lambda^t$ be the weakly continuous
Gaussian disintegration from \Cref{lem:disintegration} and define
\[
\mathfrak{J}_\lambda(t)\coloneqq
\mathfrak{J}_\lambda(\mu_\lambda^t;\mu_\lambda).
\]
For any two evidence levels $t_1,t_2$ for which the finite part is defined,
\begin{equation}
    \mathfrak{J}_\lambda(t_2)-\mathfrak{J}_\lambda(t_1)
    =
    \frac{(t_2-m_{\lambda,k})^2-(t_1-m_{\lambda,k})^2}
    {2\Sigma_{\lambda,kk}}.
\end{equation}
\begin{proof}
Let $\Delta_t\coloneqq t-m_{\lambda,k}$ and
\[
    \rho_\lambda
    \coloneqq
    \frac{\|\Sigma_{\lambda,Ik}\|_{H_{C_\lambda}}^2}{\Sigma_{\lambda,kk}}
    =
    1-\frac{S_{k,\lambda}}{\Sigma_{\lambda,kk}}.
\]
The Gaussian entropy computation in \Cref{lem:feldman_hajek} gives
\[
    \KL(\mu_{\lambda,I}^t\parallel\mu_{\lambda,I})
    =
    \frac{1}{2}\frac{\Delta_t^2}{\Sigma_{\lambda,kk}}\rho_\lambda
    -\frac{1}{2}\log(1-\rho_\lambda)-\frac{1}{2}\rho_\lambda .
\]
The conditional-residual term in \eqref{eq:unified_action} is
\[
    \frac{1}{2S_{k,\lambda}}
    \E_{\mu_\lambda^t}\!\left[
    (X_k-\mu_{k|I}(X_I))^2
    \right]
    =
    \frac{1}{2}\rho_\lambda
    +
    \frac{1}{2}\frac{\Delta_t^2}{\Sigma_{\lambda,kk}}(1-\rho_\lambda).
\]
Adding these two identities cancels the $\rho_\lambda/2$ terms and yields
\[
    \mathfrak{J}_\lambda(t)
    =
    \frac{(t-m_{\lambda,k})^2}{2\Sigma_{\lambda,kk}}
    -\frac{1}{2}\log(1-\rho_\lambda).
\]
The logarithmic term depends on $C_\lambda$ and $\lambda$, but not on $t$. Hence it
 cancels in the difference between $t_2$ and $t_1$, proving
 \eqref{eq:weak_scheme_independence}. Replacing the Gaussian Dirac mollifier in
 \Cref{prop:mollification_limit} by any centered approximate identity with finite
 second moment changes the extracted finite part only by an additive term depending on
 the kernel, $\lambda$, and $S_{k,\lambda}$, not on $t$. Therefore all evidence
 differences $\Delta\mathfrak{J}_\lambda$ are kernel-independent.
\end{proof}
\end{lemma}

\begin{proof}[Proof of Theorem~\ref{thm:finite_part_regularization}]
Put $\Delta_\lambda=x_k^*-m_{\lambda,k}$ and
$Z_\lambda=W_{\Sigma_{\lambda,Ik}}(X_I-m_{\lambda,I})$.  The Gaussian
conditional formulas give
\[
  \mathbb{E}_{\mathrm{obs}}[Z_\lambda]
  =\rho_\lambda\Delta_\lambda,
  \qquad
  \operatorname{Var}_{\mathrm{obs}}(Z_\lambda)
  =
  \Sigma_{\lambda,kk}\rho_\lambda(1-\rho_\lambda).
\]
Since $X_k=x_k^*$ under $\mu_{\mathrm{obs}}$ and
$S_{k,\lambda}=\Sigma_{\lambda,kk}(1-\rho_\lambda)$,
\[
  \frac{1}{2S_{k,\lambda}}
  \mathbb{E}_{\mu_{\mathrm{obs}}}
  \bigl[(X_k-\mu_{k|I}(X_I))^2\bigr]
  =
  \frac{\rho_\lambda}{2}
  +
  \frac{\Delta_\lambda^2}{2\Sigma_{\lambda,kk}}(1-\rho_\lambda).
\]
The calibrated identities yield the relation
$\Delta_\lambda^2 / (2\Sigma_{\lambda,kk}) = b_k^2 / \varepsilon_0$
as well as $\rho_\lambda = O(\lambda^{-4})$.  Adding
\Cref{lem:feldman_hajek} yields
\[
  \mathfrak{J}_\lambda(\mu_{\mathrm{obs}};\mu_\lambda)
  =
  \frac{b_k^2}{\varepsilon_0}+O(\lambda^{-4}).
\]
\end{proof}

If $a_{Ik}=0$ and $x_k^*=m_{\lambda,k}$, then
$\Sigma_{\lambda,Ik}=0$, the bulk marginals coincide, and the conditional
residual term in Definition~\ref{def:unified_action} is zero.  Hence
$\mathfrak{J}_\lambda(\mu_{\mathrm{obs}};\mu_\lambda)=0$ exactly for every
$\lambda>0$.

\section{Admissibility Boundary and Elliptic Benchmark}
\label{app:admissibility_boundary}

The finite-energy assumption is a boundary condition for the theory, not a
technical artifact.  The off-diagonal channel must land in the Cameron--Martin
range of the background covariance.  If it does not, the variational subtraction
need not be finite and the $O(\lambda^{-6})$ rank-one residual estimate has no
meaning.

The elliptic example of \Cref{sec:spde} makes this boundary concrete.  For
$L=-d^2/dx^2+\kappa^2$ on $L^2(0,1)$ with Dirichlet boundary conditions and
$Q_0=L^{-1}$, the Cameron--Martin space is
\[
  H_{Q_0}=H_0^1(0,1),
  \qquad
  \|h\|_{Q_0}^2
  =
  \int_0^1\bigl(|h'(x)|^2+\kappa^2|h(x)|^2\bigr)\,dx .
\]
Smooth spatially averaged observations are admissible when their representing
vectors lie in $H_0^1(0,1)$; for example, any fixed
$g\in C_c^\infty(0,1)$ has finite energy.  A point-evaluation channel is the
opposite limit.  If $g_\varepsilon$ is a standard compactly supported mollifier
concentrating at $x_0\in(0,1)$, then
\[
  \|g_\varepsilon\|_{Q_0}^2
  =
  \int_0^1\bigl(|g_\varepsilon'(x)|^2+\kappa^2|g_\varepsilon(x)|^2\bigr)\,dx
  \longrightarrow \infty
  \qquad(\varepsilon\downarrow0).
\]
Thus each fixed smoothed sensor can satisfy the finite-energy condition, but
the Dirac point-evaluation limit does not.  The shorting theorem is therefore a
theory of admissible Cameron--Martin channels, not of unsmoothed point sensors
outside the covariance-energy domain.
